\subsubsection{CFL-based Time Stepping}
\label{SEC:Time:Algorithms:CFL}

%%%%%%%%%%%%%%%%%%%%%%%%%%%%%%%%%%%%%%%%%%%%%%%%%%%%%%%%%%%%%%%%%%%%%%%%%%%%%%%%

The majority of the Runge-Kutta methods in GFR use a time step that is loosely
based on the CFL condition.
%
With \CFLmax{} being the maximum stable CFL value for a given method, the
maximum time step size, \dt, allowed for the method is
%
\[ \dt = \frac{h}{\eigmax}\:\CFLmax \]
%
where $h$ is a characteristic length scale for the cell and \eigmax{} is the
magnitude of the maximum eigenvalue for the governing equations, also sometimes
referred to as the spectral radius.
%
For traditional ERK methods, \CFLmax{} is somewhere in the range of 0.2--2.0.

%%%%%%%%%%%%%%%%%%%%%%%%%%%%%%%%%%%%%%%%%%%%%%%%%%%%%%%%%%%%%%%%%%%%%%%%%%%%%%%%

GFR uses an approximation to compute the ratio $h/\eigmax$, which means
\CFLmax{} --- which is a property of a given ERK method that can be analytically
determined --- no longer applies.
%
The \inp[nmlTimestep::CFL]{CFL} used by GFR is analogous to the true \CFLmax{}
in that it scales the GFR approximation of $h/\eigmax$ to determine \dt.
%
The GFR \inptext{CFL} is generally on the same order as \CFLmax, with stable
values ranging between 0.5--10.0 depending on the ERK method being used.

%%%%%%%%%%%%%%%%%%%%%%%%%%%%%%%%%%%%%%%%%%%%%%%%%%%%%%%%%%%%%%%%%%%%%%%%%%%%%%%%

Most ERK methods implemented in GFR have been implemented using a low-storage
algorithm that requires three registers per stage (3S*).\cite{Ketcheson_2010}
%
The low-storage 3S* algorithm for $s$ stages is shown in
\cref{ALG:Time:Algorithms:OERK}.
%
\begin{figure}[tb]
  \begin{minipage}[t]{0.49\textwidth}
    \begin{algorithm}[H]
          \centering
    \caption{Low storage implementation (3S*)}
    \label{ALG:Time:Algorithms:OERK}
    \begin{algorithmic}[1]
      \STATE $S_1 \leftarrow u^n ; S_2 \leftarrow 0 ; S_3 \leftarrow u^n$
      \FOR{$i=1:s$}
      \STATE $\tau \leftarrow t^n + c_i \Delta t$
      \STATE $S_2 \leftarrow S_2 + \delta_i S_1$
      \STATE $S_1 \leftarrow \gamma_{1,i} S_1 + 
                             \gamma_{2,i} S_2 +
                             \gamma_{3,i} S_3 + 
                      \beta_i \Delta t \bm{\mathrm{R}}\left(\tau,S_1\right)$
      \ENDFOR
      \STATE $u^{n+1} \leftarrow S_1$
    \end{algorithmic}

    \end{algorithm}
  \end{minipage}
  \hfill
  \begin{minipage}[t]{0.49\textwidth}
    \begin{algorithm}[H]
          \centering
    \caption{Embedded-pair low storage implementation (3S*$\!{}_+$)}
    \label{ALG:Time:Algorithms:EPERK}
    \begin{algorithmic}[1]
      \STATE $S_1 \leftarrow u^n ;
              S_2 \leftarrow 0 ;
              S_3 \leftarrow u^n ;
              S_4 \leftarrow u^n $
      \FOR{$i=1:s$}
      \STATE $\tau \leftarrow t^n + c_i \Delta t$
      \STATE $S_2 \leftarrow S_2 + \delta_i S_1$
      \STATE $R_i \leftarrow \bm{\mathrm{R}}\left(\tau,S_1\right)$
      \STATE $S_1 \leftarrow \gamma_{1,i} S_1 + 
                             \gamma_{2,i} S_2 +
                             \gamma_{3,i} S_3 + 
                      \beta_i \Delta t R_i$
      \STATE $S_4 \leftarrow S_4 + 
                      \hat{b}_i \Delta t R_i$
      \ENDFOR
      \STATE $u^{n+1} \leftarrow S_1$
      \STATE $\hat{u}^{n+1} \leftarrow S_4$
      \IF{$\hat{b}_{s+1} \ne 0$}
      \STATE $\hat{u}^{n+1} \leftarrow \hat{u}^{n+1} +
            \hat{b}_{s+1} \Delta t \bm{\mathrm{R}}\left(t^{n+1},u^{n+1}\right)$
      \ENDIF
    \end{algorithmic}

    \end{algorithm}
  \end{minipage}
\end{figure}
%
%
A register is essentially an array that consumes the same amount of memory as
the solution array, \thatis all the conserved variables at all the solution
points.
%
The three registers required are: 1) the actual solution array for the current
Runge-Kutta stage, $S_1$; 2) the solution array from the previous iteration,
$S_3$, and 3) an auxiliary array, $S_2$.
%
The function $\bm{\mathrm{R}}\left(t^n,u^n\right)$ represents the evaluation of
the right-hand-side residual of the system of equations, where $t^n$ and $u^n$
are respectively the time and solution at iteration $n$, and $\Delta t$ is the
change in time for the current iteration.
%
The low-storage 3S* coefficients [$\gamma_{1,i}$, $\gamma_{2,i}$,
$\gamma_{3,i}$, $\delta_i$, $\beta_i$, $c_i$] can be found in the references for
the ERK schemes.

%%%%%%%%%%%%%%%%%%%%%%%%%%%%%%%%%%%%%%%%%%%%%%%%%%%%%%%%%%%%%%%%%%%%%%%%%%%%%%%%

There are two different methods in GFR for computing an approximation to
\eigmax.
%
When
\inp[nmlTimestep::TimestepType]{$\left|\text{Timestep\_Type}\right|$}[Timestep\_Type]=1,
\eigmax{} is computed in each cell by taking the maximum value of the spectral
radius evaluated at each face point of the cell.
%
This is the recommended method to use.
%
When \inptext{$\left|\text{Timestep\_Type}\right|$}=2, \eigmax{} is computed in
each cell by taking the maximum value of the inviscid eigenvalue evaluated at
each solution point of the cell.
%
This second method is experimental and not recommended.

%%%%%%%%%%%%%%%%%%%%%%%%%%%%%%%%%%%%%%%%%%%%%%%%%%%%%%%%%%%%%%%%%%%%%%%%%%%%%%%%

\subsubsubsection{Global CFL-based Time Stepping}
\label{SEC:Time:Algorithms:CFL:Global}

Global time stepping is required for time-accurate simulations like LES.
%
When using CFL-based time stepping, a global time step is used when
$\inptext{Timestep\_Type}>0$.

%%%%%%%%%%%%%%%%%%%%%%%%%%%%%%%%%%%%%%%%%%%%%%%%%%%%%%%%%%%%%%%%%%%%%%%%%%%%%%%%

At the beginning of each solver iteration, each grid cell computes its own
unique time step size, and the minimum time step size across all grid cells
is taken to be the global time step size that all cells use to update their
solutions in time.
%
Using the same time step size for all cells ensures that the solution across the
entire computational domain is synchronized in time, and taking the minimum time
step size ensures stability for all cells.
%
This of course leads to the inefficiency in that the cell with the smallest time
step size \quoterm{slows} all the other cells down.
%
This makes global time stepping inefficient for steady-state problems.

%%%%%%%%%%%%%%%%%%%%%%%%%%%%%%%%%%%%%%%%%%%%%%%%%%%%%%%%%%%%%%%%%%%%%%%%%%%%%%%%

\subsubsubsection{Local CFL-based Time Stepping}
\label{SEC:Time:Algorithms:CFL:Local}

A local time step is used when $\inptext{Timestep\_Type}<0$.
%
For CFL-based local time stepping, each cell uses its own time step size to
advance the simulation in time.
%
At the beginning of each solver iteration, each grid cell computes its own
unique time step size that is used to update the solution within the cell at
the end of each Runge-Kutta stage.
%
Small cells with strong gradients tend to have the smallest local time step.
%
Since all cells are updated with different time step sizes, this can only be
used for steady-state simulations and cannot be used for time-accurate
simulations (\forexample LES) which require all cells to be synchronized in
time.

%%%%%%%%%%%%%%%%%%%%%%%%%%%%%%%%%%%%%%%%%%%%%%%%%%%%%%%%%%%%%%%%%%%%%%%%%%%%%%%%

